% Generated draft: requires mathematical and source review.
\documentclass[11pt]{article}
\usepackage[margin=25mm]{geometry}
\usepackage{fontspec}
\setmainfont{Latin Modern Roman}
\usepackage{amsmath,amssymb,mathtools,unicode-math}
\setmathfont{Latin Modern Math}
\usepackage{enumitem,needspace,xurl,hyperref}
\setlength{\parindent}{0pt}
\setlength{\parskip}{5pt}
\setlength{\emergencystretch}{3em}
\newcommand{\N}{\mathbb{N}}
\newcommand{\Z}{\mathbb{Z}}
\newcommand{\Q}{\mathbb{Q}}
\newcommand{\R}{\mathbb{R}}
\newcommand{\C}{\mathbb{C}}
\newcommand{\F}{\mathbb{F}}
\newcommand{\A}{\mathbb{A}}
\newcommand{\PP}{\mathbb{P}}
\newcommand{\E}{\mathbb{E}}
\newcommand{\eps}{\varepsilon}
\newcommand{\dd}{\,\mathrm{d}}
\begin{document}
\section*{UnsolvedMath ID: 700636}
\textbf{Source problem code:} catalogue entry
\subsection{Definitions and mathematical statement}
Let \(R\) be an associative ring, not assumed commutative or to have a multiplicative identity. Interpret \emph{radical} in the standard Jacobson radical sense: every \(a\in R\) is quasi-regular, meaning that there exists \(b\in R\) such that
\[
a+b+ab=0=b+a+ba.
\]
The source does not explicitly specify this convention. Define the adjoint group \(R^{\circ}\) to have underlying set \(R\), identity \(0\), and multiplication
\[
a\circ b=a+b+ab.
\]
Define the associated Lie ring \(R^{-}\) to have the additive group of \(R\) and bracket \([a,b]=ab-ba\). A Lie ideal of \(R^{-}\) means an additive subgroup \(I\leq (R,+)\) satisfying \([r,i]\in I\) for all \(r\in R\) and \(i\in I\).

Characterize all radical associative rings \(R\) for which
\[
\bigl\{N\subseteq R:N\trianglelefteq R^{\circ}\bigr\}
=
\bigl\{I\subseteq R:I\leq(R,+),\ [r,i]\in I\text{ for every }r\in R,\ i\in I\bigr\}.
\]
The equality is equality of families of subsets of the common underlying set \(R\). Thus it requires both that every normal subgroup of the adjoint group be a Lie ideal and that every Lie ideal be a normal subgroup of the adjoint group. No finiteness, characteristic, or base-field restriction is supplied.
\subsection{Short English statement}
Characterize the radical associative rings whose adjoint-group normal subgroups are exactly the ideals of their associated Lie rings, viewed as subsets of the same ring. Here radical is interpreted as Jacobson radical, and Lie ideals are additive subgroups stable under brackets with every ring element.
\subsection{Sources}
\begin{itemize}
\item Local UnsolvedMath registry, permanent ID \(700636\), exact source problem code \texttt{catalogue entry}, title \emph{Kourovka Notebook Problem 10.19}; supplied upstream revision \texttt{d371ee917ac9647fbb06cd6fc529deec6f3dda82}.
\item \emph{The Kourovka Notebook (Unsolved Problems in Group Theory)}, original problem 10.19, 10th issue (1986); proposed by V. M. Levchuk.
\item The supplied record identifies the 21st issue edition (2026), arXiv:1401.0300v48. Its collection description names E. I. Khukhro and V. D. Mazurov as editors. These bibliographic details are supplied attribution and have not been independently verified.
\item The registry labels this problem \emph{open}. This is a supplied status, not an independently verified literature assessment.
\end{itemize}
\end{document}
